\section{Boolean analysis and expander walks}

\begin{result}
\subsection{Hypercontractivity on the Boolean cube}
For a real-valued function on a finite Boolean cube with uniform probability
measure, $1\le p\le2\le q<\infty$, and
$0\le\rho\le\sqrt{(p-1)/(q-1)}$, the noise operator satisfies
\[
\|T_\rho f\|_q\le\|f\|_p.
\]
This includes the $(2,4)$ inequality at $\rho=1/\sqrt3$. The indexed theorem
covers the range $p\le2\le q$ and finite real exponents; it is not a claim of all
hypercontractive parameter regimes.\citep{Bonami1970,ODonnell2014}
\decl{Complexity/Foundations/BooleanAnalysis/Hypercontractivity.lean}{202}{Cslib.BitString.lpNorm_noiseOp_le_lpNorm}
\end{result}

\begin{result}
\subsection{The Kahn--Kalai--Linial theorem and Talagrand's inequality}
Let $f$ be a Boolean function on a finite cube with $n\ge1$ coordinates, read as
the sign $F=(-1)^{f}$, with influences $\mathrm{Inf}_i[F]$ (the probability that
flipping bit $i$ changes $f$), total influence $I[F]=\sum_i\mathrm{Inf}_i[F]$, and
variance $\mathrm{Var}[F]$, all for the uniform measure. Then some coordinate has
\[
\mathrm{Inf}_i[F]\ge\frac{\mathrm{Var}[F]\,\log n}{12\,n}.
\]
This follows from the edge-isoperimetric form: if every influence is at most
$\delta$, then $\mathrm{Var}[F]\log(1/\delta)\le 6\,I[F]$. If every influence is
less than one, $\mathrm{Var}[F]\le6\sum_i\mathrm{Inf}_i[F]/\log(1/\mathrm{Inf}_i[F])$
(Talagrand), and for every $f$,
$\mathrm{Var}[F]\le7\sum_i\mathrm{Inf}_i[F]/(1+\log(1/\mathrm{Inf}_i[F]))$; terms with
$\mathrm{Inf}_i[F]=0$ contribute $0$. Logarithms are natural; constants are those of the
formal proof, not optimized. The proof applies the library's $(4/3,2)$-hypercontractive
inequality to the derivatives $D_iF$.\citep{KahnKalaiLinial1988,Talagrand1994,ODonnell2014}
\declat{f355a73f091795556fe6a70dc419d568d3a24df9}{Complexity/Foundations/BooleanAnalysis/KahnKalaiLinial.lean}{310}{Cslib.BitString.exists_expect_sub_expect_sq_mul_log_card_div_le_influence}
\declat{f355a73f091795556fe6a70dc419d568d3a24df9}{Complexity/Foundations/BooleanAnalysis/KahnKalaiLinial.lean}{294}{Cslib.BitString.expect_sub_expect_sq_mul_log_inv_le_six_mul_totalInfluence}
\declat{f355a73f091795556fe6a70dc419d568d3a24df9}{Complexity/Foundations/BooleanAnalysis/KahnKalaiLinial.lean}{371}{Cslib.BitString.expect_sub_expect_sq_le_six_mul_sum_influence_div_log}
\declat{f355a73f091795556fe6a70dc419d568d3a24df9}{Complexity/Foundations/BooleanAnalysis/KahnKalaiLinial.lean}{390}{Cslib.BitString.expect_sub_expect_sq_le_seven_mul_sum_influence_div_one_add_log}
\end{result}

\begin{result}
\subsection{Friedgut's junta theorem}
Let $f$ be a Boolean function on a finite cube whose sign $F=(-1)^f$ has total
influence $I[F]\le K$, and let $\varepsilon>0$. Then there are a set $J$ of at most
$\frac{4K^3}{9\varepsilon^2}\,81^{K/\varepsilon}$ coordinates and a Boolean function
$g$ depending only on the coordinates in $J$ with $\Pr_x[f(x)\ne g(x)]\le\varepsilon$;
for $0<\varepsilon\le1$ one may take $|J|\le3\cdot2^{8K/\varepsilon}$. In the proof, $J$ is
the set of coordinates of influence above an explicit threshold and $g$ is the sign of
the average of $F$ over the other coordinates. The low-degree part of O'Donnell's
spectral concentration statement (sets of size at most $I[F]/\varepsilon$) is not
formalized.\citep{Friedgut1998,ODonnell2014}
\declat{f355a73f091795556fe6a70dc419d568d3a24df9}{Complexity/Foundations/BooleanAnalysis/FriedgutJunta.lean}{217}{Cslib.BitString.exists_dependsOn_hammingDist_div_le_of_totalInfluence_le}
\declat{f355a73f091795556fe6a70dc419d568d3a24df9}{Complexity/Foundations/BooleanAnalysis/FriedgutJunta.lean}{285}{Cslib.BitString.exists_card_le_three_mul_two_rpow_dependsOn_hammingDist_div_le}
\end{result}


\begin{result}
\subsection{Expander-walk hitting bound}
Let a finite undirected regular graph be represented by a rotation map with
nonempty vertex and port sets, and let $\lambda$ be its nonconstant spectral
operator-norm bound. Start a walk at a uniform vertex and choose each port
uniformly. If each set $B_i$ has density at most $0\le\mu\le1$, then
\[
\Pr[v_i\in B_i\text{ for every }0\le i\le t]
 \le \mu\bigl(\mu+\lambda(1-\mu)\bigr)^t.
\]
The sets may vary with time. This is a finite probability bound; machine-level
error reduction and extractor evaluation costs are separate obligations.
\citep{AjtaiKomlosSzemeredi1987,Kahale1995,Vadhan2012}
\decl{Complexity/Foundations/Combinatorics/Expander/Hitting.lean}{207}{Cslib.RotationMap.dens_filter_forall_walkAt_mem_le}
\end{result}

\begin{result}
\subsection{The Nisan--Szegedy bound on relevant coordinates}
If a Boolean function on a finite cube, read as the sign $(-1)^{f(x)}$,
has no nonzero Fourier coefficients on sets of more than $d$ coordinates,
then at most $d\,2^{d-1}$ coordinates are relevant: flipping an irrelevant
coordinate never changes the value. For $d=0$ the function has no relevant
coordinates; the formal natural exponent uses truncated subtraction.
The Fourier coefficients use the uniform measure on the cube.
\citep{NisanSzegedy1994,ODonnell2014}
\decl{Complexity/Foundations/BooleanAnalysis/RelevantCoordinates.lean}{103}{Cslib.BitString.card_filter_exists_update_ne_le}
\end{result}

\begin{result}
\subsection{Fourier bounds for decision trees}
Suppose a deterministic decision tree has structural depth $d$ and $L$ leaves.
For any real-valued readout of its output bounded in absolute value by one,
the sum of absolute Fourier coefficients is at most $L\le2^d$.
For a Boolean readout, written as signs, each coefficient is an integer multiple
of $2^{-\max(d-1,0)}$, and at most $L\,2^{\max(d-1,0)}$ coefficients are nonzero.
Repeated queries and unreachable branches are allowed; structural size only
weakens these upper bounds. This is a spectral statement, not a formalized
learning algorithm.\citep{KushilevitzMansour1993,ODonnell2014}
\decl{Complexity/Computability/DecisionTree/SpectralNorm.lean}{72}{Cslib.DecisionTree.sum_abs_fourierCoeff_le_leafCount}
\decl{Complexity/Computability/DecisionTree/Granularity.lean}{84}{Cslib.DecisionTree.exists_fourierCoeff_signVal_eq_intCast_div_two_pow_depth_sub_one}
\decl{Complexity/Computability/DecisionTree/SpectralNorm.lean}{113}{Cslib.DecisionTree.card_filter_fourierCoeff_signVal_ne_zero_le}
\end{result}

\begin{result}
\subsection{A degree-dependent obstruction to perfect expansion}
For an undirected $D$-regular multigraph on $N\ge2$ vertices, with $D\ge1$,
let $\lambda$ be the norm of its normalized walk operator on mean-zero functions.
Then
\[
\frac{N-D}{D(N-1)}\le\lambda^2.
\]
In particular $\lambda\ge1/\sqrt{2D}$ if $N\ge2D$.
The rotation-map model permits loops and parallel edges. The complete graph
with one loop at each vertex has $D=N$ and $\lambda=0$, attaining the first
bound. This finite trace bound is not the asymptotic Alon--Boppana theorem.
\citep{HooryLinialWigderson2006}
\decl{Complexity/Foundations/Combinatorics/Expander/LowerBound.lean}{91}{Cslib.RotationMap.div_le_sq_spectralExpansion}
\decl{Complexity/Foundations/Combinatorics/Expander/LowerBound.lean}{195}{Cslib.RotationMap.one_div_sqrt_two_mul_le_spectralExpansion}
\decl{Complexity/Foundations/Combinatorics/Expander/LowerBound.lean}{77}{Cslib.RotationMap.spectralExpansion_complete}
\end{result}
